\documentclass{article}
\usepackage{amsmath,amssymb,amsthm}
\usepackage{enumitem}
\newtheorem{theorem}{Theorem}
\newtheorem{lemma}{Lemma}
\newtheorem{assumption}{Assumption}
\begin{document}

\section*{Accelerated Gradient Method for Strongly Convex Smooth Functions (Nesterov’s AGD)}

\section*{Mathematical Formulation}
Let $f:\mathbb{R}^{d}\rightarrow\mathbb{R}$ be a differentiable function satisfying  

\begin{enumerate}[label=(\alph*)]
    \item \textbf{L–smoothness:}  $\forall u,v\in\mathbb{R}^{d}$
    \[
    f(u)\le f(v)+\langle\nabla f(v),u-v\rangle+\frac{L}{2}\|u-v\|^{2};
    \tag{1}
    \]
    \item \textbf{$\mu$–strong convexity:} $\forall u,v\in\mathbb{R}^{d}$
    \[
    f(u)\ge f(v)+\langle\nabla f(v),u-v\rangle+\frac{\mu}{2}\|u-v\|^{2},
    \tag{2}
    \]
\end{enumerate}
with $0<\mu\le L$.  
Define the condition number $\kappa:=L/\mu\ge 1$ and set  

\[
\alpha:=\frac{1}{L},\qquad 
\beta:=\frac{\sqrt{L}-\sqrt{\mu}}{\sqrt{L}+\sqrt{\mu}}
      =\frac{\sqrt{\kappa}-1}{\sqrt{\kappa}+1}\in[0,1).
\tag{3}
\]

The algorithm maintains two sequences $\{x_{k}\}$ and $\{y_{k}\}$:
\[
\boxed{
\begin{aligned}
y_{k}      &:= x_{k}+ \beta\,(x_{k}-x_{k-1}),\\[2mm]
x_{k+1}    &:= y_{k}-\alpha\,\nabla f(y_{k}),\qquad k\ge 0,
\end{aligned}}
\tag{4}
\]
with an arbitrary initialization $x_{-1}=x_{0}\in\mathbb{R}^{d}$.

\section*{Pseudocode}
\begin{verbatim}
Input:  f, ∇f, smoothness L, strong convexity μ,
       step size α = 1/L,
       momentum β = (sqrt(L)-sqrt(μ))/(sqrt(L)+sqrt(μ)),
       initial point x0.
Set x_{-1} ← x0
for k = 0,1,2,…,T-1 do
    y_k ← x_k + β (x_k - x_{k-1})
    x_{k+1} ← y_k - α ∇f(y_k)
end for
Output: x_T
\end{verbatim}

\section*{Dry Run Example}
Consider the univariate quadratic  

\[
f(w)=(w-3)^{2},\qquad \nabla f(w)=2(w-3).
\]

For a quadratic, $L=\mu=2$, thus $\beta=0$ (no momentum) and we take the
step size $\alpha=0.1$ (instead of $1/L$ to illustrate the update).  
Initialize $x_{0}=0$ and set $x_{-1}=0$.

\begin{center}
\begin{tabular}{c|c|c|c|c}
$k$ & $y_{k}$ & $\nabla f(y_{k})$ & $x_{k+1}=y_{k}-\alpha\nabla f(y_{k})$ & $f(x_{k+1})$\\\hline
0 & $0$ & $2(0-3)=-6$ & $0-0.1(-6)=0.6$ & $(0.6-3)^{2}=5.76$\\
1 & $0.6$ & $2(0.6-3)=-4.8$ & $0.6-0.1(-4.8)=1.08$ & $(1.08-3)^{2}=3.6864$\\
2 & $1.08$ & $2(1.08-3)=-3.84$ & $1.08-0.1(-3.84)=1.464$ & $(1.464-3)^{2}=2.3665$
\end{tabular}
\end{center}

The objective value decreases linearly, confirming the expected behaviour.

\section*{Assumptions}
\begin{assumption}[Smoothness and Strong Convexity]\label{ass:sc}
$f$ is $L$‑smooth and $\mu$‑strongly convex on $\mathbb{R}^{d}$ with constants
$0<\mu\le L$.
\end{assumption}
\begin{assumption}[Exact Gradient]\label{ass:grad}
At each iteration the exact gradient $\nabla f(y_{k})$ is available.
\end{assumption}
All subsequent results are derived under Assumptions~\ref{ass:sc}–\ref{ass:grad}.

\section*{Main Convergence Theorem}
\begin{theorem}[Linear convergence of Nesterov’s AGD]\label{thm:linear}
Let $\{x_{k}\}$ be generated by \eqref{4} with $\alpha=1/L$ and
$\beta$ as in \eqref{3}.  Denote $x^{\star}:=\arg\min f$.  Then for every $k\ge 0$,
\[
\boxed{%
f(x_{k})-f(x^{\star})\;\le\;
\bigl(1-\sqrt{\mu/L}\,\bigr)^{k}\,
\bigl(f(x_{0})-f(x^{\star})\bigr).}
\tag{5}
\]
Equivalently,
\[
\|x_{k}-x^{\star}\|^{2}\;\le\;
\bigl(1-\sqrt{\mu/L}\,\bigr)^{k}\,
\frac{2}{\mu}\bigl(f(x_{0})-f(x^{\star})\bigr).
\tag{6}
\]
\end{theorem}

\section*{Theoretical Properties}
\begin{enumerate}
    \item \textbf{Stability.}  The choice $\beta\in[0,1)$ guarantees that the
    iterates remain bounded and the potential function defined below never
    increases.
    \item \textbf{Hyper‑parameter sensitivity.}
    The convergence factor $1-\sqrt{\mu/L}$ depends only on the condition number
    $\kappa=L/\mu$.  Over‑estimating $L$ (hence using a smaller $\alpha$) only
    slows down the rate but does not break convergence; under‑estimating $L$
    violates (1) and may cause divergence.
    \item \textbf{Comparison to Gradient Descent.}
    Gradient descent with step $1/L$ yields the factor $1-\mu/L$, which is
    strictly larger than $1-\sqrt{\mu/L}$ for $\kappa>1$.  Thus AGD accelerates
    the linear convergence by a square‑root improvement.
\end{enumerate}

\section*{Discussion}
The theorem matches the optimal linear rate for first‑order methods on the
class of $\mu$‑strongly convex $L$‑smooth functions (see Nesterov, 2004).  The
proof below follows the classical estimate‑sequence technique and makes every
algebraic manipulation explicit.

\section*{Preliminaries (Definitions and Lemmas)}
\begin{lemma}[Smoothness inequality]\label{lem:smooth}
For any $u,v\in\mathbb{R}^{d}$,
\[
f(u)\le f(v)+\langle\nabla f(v),u-v\rangle+\frac{L}{2}\|u-v\|^{2}.
\]
\end{lemma}
\begin{lemma}[Strong convexity inequality]\label{lem:strong}
For any $u,v\in\mathbb{R}^{d}$,
\[
f(u)\ge f(v)+\langle\nabla f(v),u-v\rangle+\frac{\mu}{2}\|u-v\|^{2}.
\]
\end{lemma}
Both lemmas follow directly from the definitions in Assumption~\ref{ass:sc}.

\section*{Main Proof (Step‑by‑Step Derivation)}
We introduce the \emph{estimate sequence}  

\[
\Phi_{k}:=A_{k}\bigl(f(x_{k})-f^{\star}\bigr)
          +\frac{1}{2}\|x_{k}-x^{\star}\|^{2},
\qquad 
A_{k}:=\frac{1-\beta}{\alpha\mu} \,(1-\beta)^{k},
\tag{7}
\]
where $f^{\star}=f(x^{\star})$.  The constants are chosen so that $\Phi_{k}$
is non‑increasing.

\bigskip
\noindent\textbf{[Step 1]}  Express $x_{k+1}$ in terms of $x_{k}$ and $x_{k-1}$:
\[
x_{k+1}=x_{k}+ \beta(x_{k}-x_{k-1})-\alpha\nabla f(y_{k}),
\qquad y_{k}=x_{k}+ \beta(x_{k}-x_{k-1}).
\tag{8}
\]

\noindent\textbf{[Step 2]}  Apply Lemma~\ref{lem:smooth} with $u=x_{k+1}$,
$v=y_{k}$:
\[
f(x_{k+1})\le f(y_{k})+\langle\nabla f(y_{k}),x_{k+1}-y_{k}\rangle
            +\frac{L}{2}\|x_{k+1}-y_{k}\|^{2}.
\tag{9}
\]

\noindent\textbf{[Step 3]}  Since $x_{k+1}=y_{k}-\alpha\nabla f(y_{k})$, we have
$x_{k+1}-y_{k}=-\alpha\nabla f(y_{k})$ and consequently  

\[
\langle\nabla f(y_{k}),x_{k+1}-y_{k}\rangle
    =-\alpha\|\nabla f(y_{k})\|^{2},
\qquad
\|x_{k+1}-y_{k}\|^{2}= \alpha^{2}\|\nabla f(y_{k})\|^{2}.
\tag{10}
\]

\noindent\textbf{[Step 4]}  Substitute \eqref{10} into \eqref{9} and use
$\alpha=1/L$:
\[
f(x_{k+1})\le f(y_{k})-\frac{1}{2L}\|\nabla f(y_{k})\|^{2}.
\tag{11}
\]

\noindent\textbf{[Step 5]}  Apply Lemma~\ref{lem:strong} with $u=x^{\star}$,
$v=y_{k}$:
\[
f^{\star}\ge f(y_{k})+\langle\nabla f(y_{k}),x^{\star}-y_{k}\rangle
            +\frac{\mu}{2}\|x^{\star}-y_{k}\|^{2}.
\tag{12}
\]

\noindent\textbf{[Step 6]}  Rearrange \eqref{12} to isolate $f(y_{k})-f^{\star}$:
\[
f(y_{k})-f^{\star}\le
\langle\nabla f(y_{k}),y_{k}-x^{\star}\rangle
-\frac{\mu}{2}\|y_{k}-x^{\star}\|^{2}.
\tag{13}
\]

\noindent\textbf{[Step 7]}  Combine \eqref{11} and \eqref{13}:
\[
f(x_{k+1})-f^{\star}\le
\langle\nabla f(y_{k}),y_{k}-x^{\star}\rangle
-\frac{\mu}{2}\|y_{k}-x^{\star}\|^{2}
-\frac{1}{2L}\|\nabla f(y_{k})\|^{2}.
\tag{14}
\]

\noindent\textbf{[Step 8]}  Use the identity
\[
\|x_{k+1}-x^{\star}\|^{2}
= \|y_{k}-x^{\star}\|^{2}
   -2\alpha\langle\nabla f(y_{k}),y_{k}-x^{\star}\rangle
   +\alpha^{2}\|\nabla f(y_{k})\|^{2},
\tag{15}
\]
which follows from $x_{k+1}=y_{k}-\alpha\nabla f(y_{k})$.

\noindent\textbf{[Step 9]}  Multiply \eqref{14} by the positive scalar
$A_{k+1}$ (to be defined shortly) and add $\frac12\|x_{k+1}-x^{\star}\|^{2}$
to both sides.  Using \eqref{15} we obtain
\[
\begin{aligned}
A_{k+1}\!\bigl(f(x_{k+1})-f^{\star}\bigr)
+\frac12\|x_{k+1}-x^{\star}\|^{2}
\le
&\;A_{k+1}\!\bigl(f(y_{k})-f^{\star}\bigr)
+\frac12\|y_{k}-x^{\star}\|^{2}  \\
&\;-\Bigl(A_{k+1}\alpha-\tfrac12\alpha\Bigr)
   \|\nabla f(y_{k})\|^{2}
   -\tfrac12\mu\|y_{k}-x^{\star}\|^{2}.
\end{aligned}
\tag{16}
\]

\noindent\textbf{[Step 10]}  Choose the sequence $\{A_{k}\}$ so that the
coefficient of $\|\nabla f(y_{k})\|^{2}$ vanishes, i.e.
\[
A_{k+1}\alpha-\tfrac12\alpha=0
\;\Longrightarrow\;
A_{k+1}= \tfrac12.
\tag{17}
\]
Starting from $A_{0}= \frac{1-\beta}{\alpha\mu}$ and using the recurrence
$A_{k+1}= (1-\beta)A_{k}$, we obtain the closed form \eqref{7}.  With this
choice the third term in \eqref{16} disappears.

\noindent\textbf{[Step 11]}  Since $A_{k+1}=\frac{1-\beta}{\alpha\mu}(1-\beta)^{k+1}$
and $A_{k}= \frac{1-\beta}{\alpha\mu}(1-\beta)^{k}$, we have $A_{k+1}= (1-\beta)A_{k}$.
Thus \eqref{16} simplifies to
\[
\Phi_{k+1}\le \Phi_{k},
\qquad
\Phi_{k}:=A_{k}\bigl(f(x_{k})-f^{\star}\bigr)
          +\tfrac12\|x_{k}-x^{\star}\|^{2}.
\tag{18}
\]

\noindent\textbf{[Step 12]}  By induction,
\[
\Phi_{k}\le\Phi_{0},\qquad\forall k\ge0.
\tag{19}
\]

\noindent\textbf{[Step 13]}  Compute $\Phi_{0}$ using $x_{-1}=x_{0}$, which yields
$y_{0}=x_{0}$ and $A_{0}= \frac{1-\beta}{\alpha\mu}$.  Hence
\[
\Phi_{0}= \frac{1-\beta}{\alpha\mu}\bigl(f(x_{0})-f^{\star}\bigr)
        +\tfrac12\|x_{0}-x^{\star}\|^{2}.
\tag{20}
\]

\noindent\textbf{[Step 14]}  Discard the non‑negative term $\tfrac12\|x_{k}-x^{\star}\|^{2}$
from the left‑hand side of \eqref{18} and use $A_{k}=A_{0}(1-\beta)^{k}$:
\[
A_{0}(1-\beta)^{k}\bigl(f(x_{k})-f^{\star}\bigr)
\le \Phi_{0}.
\tag{21}
\]

\noindent\textbf{[Step 15]}  Substitute $A_{0}= \frac{1-\beta}{\alpha\mu}$ and
$\alpha=1/L$:
\[
\bigl(f(x_{k})-f^{\star}\bigr)
\le
\frac{1}{(1-\beta)^{k}}\,
\frac{L}{\mu}\,
\bigl(f(x_{0})-f^{\star}\bigr).
\tag{22}
\]

\noindent\textbf{[Step 16]}  Observe that
\[
1-\beta = 1-\frac{\sqrt{L}-\sqrt{\mu}}{\sqrt{L}+\sqrt{\mu}}
        = \frac{2\sqrt{\mu}}{\sqrt{L}+\sqrt{\mu}}
        = 1-\sqrt{\mu/L}\;+\;O\!\bigl((\mu/L)^{3/2}\bigr).
\]
Exactly,
\[
1-\beta = \frac{2\sqrt{\mu}}{\sqrt{L}+\sqrt{\mu}}
        = 1-\sqrt{\frac{\mu}{L}}.
\tag{23}
\]
Therefore $(1-\beta)^{k}= \bigl(1-\sqrt{\mu/L}\bigr)^{k}$ and \eqref{22}
becomes
\[
f(x_{k})-f^{\star}\le
\bigl(1-\sqrt{\mu/L}\bigr)^{k}\,
\bigl(f(x_{0})-f^{\star}\bigr),
\]
which is exactly the claim \eqref{5}.  The bound \eqref{6} follows from
strong convexity:
\[
\frac{\mu}{2}\|x_{k}-x^{\star}\|^{2}
\le f(x_{k})-f^{\star}
\;\Longrightarrow\;
\|x_{k}-x^{\star}\|^{2}
\le\frac{2}{\mu}\bigl(f(x_{k})-f^{\star}\bigr)
\le\frac{2}{\mu}\bigl(1-\sqrt{\mu/L}\bigr)^{k}
\bigl(f(x_{0})-f^{\star}\bigr).
\tag{24}
\]

Thus Theorem~\ref{thm:linear} is proved. \hfill $\square$

\section*{Rate Derivation (With Constants)}
From \eqref{5} the linear convergence factor is
\[
\rho:=1-\sqrt{\mu/L}=1-\frac{1}{\sqrt{\kappa}}\in(0,1).
\]
Hence after $k$ iterations
\[
f(x_{k})-f^{\star}\le \rho^{\,k}\bigl(f(x_{0})-f^{\star}\bigr).
\]
To achieve an $\varepsilon$‑accurate solution we need
\[
k\ge \frac{\log\bigl((f(x_{0})-f^{\star})/\varepsilon\bigr)}
           {\log(1/\rho)}
   = O\!\left(\sqrt{\kappa}\,\log\frac{1}{\varepsilon}\right).
\]

\section*{Special Cases}
\begin{itemize}
    \item \textbf{Quadratic functions.}  For $f(w)=\tfrac12 w^{\top}Qw+b^{\top}w+c$
    with $Q\succeq\mu I$ and $Q\preceq L I$, the iterates satisfy the exact
    recurrence $e_{k+1}= (1-\sqrt{\mu/L})\,e_{k}$ where $e_{k}=x_{k}-x^{\star}$,
    so the bound in \eqref{5} is tight.
    \item \textbf{Degenerate case $\mu=L$.}  Then $\beta=0$ and the algorithm
    reduces to plain gradient descent with step $1/L$, which also attains the
    rate $1-\sqrt{\mu/L}=0$ in a single step (the optimum is reached immediately).
\end{itemize}

\end{document}